\chapter{Extinction and Component Topology}

This chapter proves finite extinction for the controlled surgery flow.  The
argument first eliminates components with nontrivial second homotopy, then
uses a loop-space width attached to a nonzero third-homotopy class.

\section{Two-Sphere Width Control}

For a $C^1$ map $f\colon S^2\to X$ into a Riemannian manifold, its area is
the integral of the two-dimensional Jacobian, counted with multiplicity.

\begin{definition}[Least nontrivial two-sphere area]
\label{def:two-sphere-width}
\dcref{MT:ch18:18.10}
\notready
Let $(X,g)$ be compact with $\pi_2(X)\ne0$.  Its \emph{two-sphere width} is
\[
 W_2(X,g)=\inf\{\operatorname{Area}_g(f):f\colon S^2\to X\text{ is $C^1$ and not
 null-homotopic}\}.
\]
For a component path $\mathcal X$ in a Ricci flow with surgery, write
\[
 W_2(t)=W_2(\mathcal X(t),g(t))
\]
whenever
$\pi_2(\mathcal X(t))\ne0$.
\end{definition}

\begin{theorem}[Realization of the two-sphere width]
\label{thm:two-sphere-width-realization}
\dcref{MT:ch18:18.10}
\uses{def:two-sphere-width}
\notready
Assume the Sacks--Uhlenbeck minimization and regularity theorem for the
nonzero $\pi_2$-class in this compact three-manifold.  If $(X,g)$ is compact
and $\pi_2(X)\ne0$, then $W_2(X,g)>0$ and there is a
non-null-homotopic conformal harmonic map $f\colon S^2\to X$, a branched
minimal immersion away from finitely many points, such that
\[
 \operatorname{Area}_g(f)=W_2(X,g).
\]
\end{theorem}

\begin{lemma}[Area variation of a minimal two-sphere]
\label{lem:minimal-two-sphere-area-variation}
\dcref{MT:ch18:18.12}
\notready
Let $g(t)$ be a Ricci flow on a three-manifold near $t_0$, and let
$f\colon S^2\to(M,g(t_0))$ be a branched minimal immersion.  If
$R_{\min}(t_0)=\min_M R_{g(t_0)}$, then
\[
 \left.\frac d{dt}\right|_{t=t_0}\operatorname{Area}_{g(t)}(f)
 \le -4\pi-\frac12R_{\min}(t_0)\operatorname{Area}_{g(t_0)}(f).
\]
\end{lemma}

\begin{proof}
Away from the finite branch set, differentiation of the induced area form
under $\partial_tg=-2\operatorname{Ric}$ gives
\[
 \frac d{dt}\operatorname{Area}_{g(t)}(f)
 =-\int_{S^2}\bigl(R-\operatorname{Ric}(\nu,\nu)\bigr)\,da.
\]
For a minimal surface, the Gauss equation rewrites the integrand as
\[
 -\int_{S^2}K_f\,da-
 \frac12\int_{S^2}(|A|^2+R)\,da.
\]
If the branch orders are $b_1,\ldots,b_q\ge1$, branched Gauss--Bonnet gives
$\int K_f\,da=4\pi+2\pi\sum_jb_j$.  The branch and second-fundamental-form
terms therefore have nonpositive sign.  Finally
$\int R\,da\ge R_{\min}(t_0)\operatorname{Area}_{g(t_0)}(f)$, which proves the claim.
\end{proof}

\begin{lemma}[Fixed-map area distortion]
\label{lem:fixed-sphere-area-distortion}
\dcref{MT:ch18:18.13}
\notready
Let $g(t)$ be a Ricci flow on a compact manifold for
$t\in[t_0,t_1]$, suppose $|\operatorname{Ric}_{g(t)}|\le D$, and let
$f\colon S^2\to M$ be $C^1$.  Then, for $s,t\in[t_0,t_1]$,
\[
 e^{-4D|t-s|}\operatorname{Area}_{g(s)}(f)
 \le \operatorname{Area}_{g(t)}(f)
 \le e^{4D|t-s|}\operatorname{Area}_{g(s)}(f).
\]
\end{lemma}

\begin{proof}
Put $A(t)=\operatorname{Area}_{g(t)}(f)$.  At almost every regular point of $f$, the
variation of the two-dimensional Jacobian is one half of the trace of
$\partial_tg=-2\operatorname{Ric}$ on the image plane.  The source norm convention gives
$|A'(t)|\le4DA(t)$.  Thus
$e^{-4D(t-s)}A(t)$ is nonincreasing for $t\ge s$, which proves the upper
bound.  Interchanging $s$ and $t$ proves the lower bound.  Approximation by
smooth maps and dominated convergence remove the regularity restriction.
\end{proof}

\begin{lemma}[Continuity of the two-sphere width at regular times]
\label{lem:two-sphere-width-regular-continuity}
\dcref{MT:ch18:18.11,MT:ch18:18.13}
\uses{def:two-sphere-width,thm:two-sphere-width-realization,
  lem:fixed-sphere-area-distortion}
\notready
On a closed surgery-free interval on which the component is identified by
the Ricci-flow time vector field and has nonzero $\pi_2$, the function $W_2$
is continuous.
\end{lemma}

\begin{proof}
Fix $t$ and let $t_i\to t$.  A minimizing sphere at $t$, supplied by
\ref{thm:two-sphere-width-realization}, represents the same nonzero class at
nearby times.  The upper estimate in
\ref{lem:fixed-sphere-area-distortion} gives
$\limsup_iW_2(t_i)\le W_2(t)$.  Conversely, choose a minimizing sphere at
each $t_i$ and view it at time $t$.  The reverse estimate, whose constant is
uniform for $t_i$ near $t$, gives
$W_2(t)\le\liminf_iW_2(t_i)$.  The two inequalities prove continuity.
\end{proof}

\begin{theorem}[Component map for a trivial surgery]
\label{thm:trivial-surgery-component-map}
\dcref{MT:ch18:18.18,MT:ch18:18.22,MT:ch15:15.12}
\uses{thm:separating-surgery-map}
\notready
Let $T$ be a surgery time, let $X(t)$ be a presurgery component for $t<T$
close to $T$, and suppose every cutting sphere in this component is null
homotopic.  Assume the connected-sum analysis supplies, for each selected
postsurgery component $Y$, a pinch map with the stated $\pi_2$ uniqueness and
homotopy-equivalence properties, and a nonzero-degree map on $\pi_3$ whenever
$\pi_3(X(t))\cong\mathbb Z$.  Then for every $\eta>0$ there is, for all
$t<T$ sufficiently close to $T$,
\[
 F_t\colon X(t)\longrightarrow Y,
 \qquad \operatorname{Lip}(F_t)\le1+\eta,
\]
homotopic to the corresponding connected-sum pinch map.  The component with
nonzero $\pi_2$, when it exists, is unique and its pinch map is a homotopy
equivalence.  If $\pi_2(X(t))=0$, then $\pi_2(Y)=0$.  If in addition
$\pi_3(X(t))\cong\mathbb Z$, then $\pi_3(Y)\cong\mathbb Z$ and $(F_t)_*$ is
multiplication by a nonzero integer; in particular it is injective.
\end{theorem}

\begin{proof}
The separating-surgery map \ref{thm:separating-surgery-map} supplies the
$(1+\eta)$-Lipschitz representative.  The connected-sum hypothesis in the
statement supplies its pinch-map homotopy and homotopy-group conclusions.
\end{proof}

\begin{lemma}[Two-sphere width across a trivial surgery]
\label{lem:two-sphere-width-surgery-control}
\dcref{MT:ch18:18.11,MT:ch15:15.12}
\uses{def:two-sphere-width,thm:trivial-surgery-component-map}
\notready
Let $T$ be a surgery time after which the surgery spheres in a continuing
component are null homotopic, and suppose the selected postsurgery component
has nonzero $\pi_2$.  Then
\[
 W_2(T)\le\liminf_{t\uparrow T}W_2(t).
\]
\end{lemma}

\begin{proof}
By \ref{thm:trivial-surgery-component-map}, the postsurgery component with
nonzero $\pi_2$ is unique and its pinch map is a homotopy equivalence.  For
every $\eta>0$ and all
$t<T$ sufficiently close to $T$, the pinch equivalence is represented by a
map $F_t\colon\mathcal X(t)\to\mathcal X(T)$ with Lipschitz constant at most
$1+\eta$.  Composition with $F_t$ preserves nontriviality of
a sphere class and multiplies area by at most $(1+\eta)^2$.  Hence
$W_2(T)\le(1+\eta)^2W_2(t)$.  Take a sequence realizing the left liminf and
then let $\eta\downarrow0$.
\end{proof}

\begin{lemma}[Forward inequality for the two-sphere width]
\label{lem:two-sphere-width-forward-inequality}
\dcref{MT:ch18:18.11}
\uses{def:two-sphere-width,thm:two-sphere-width-realization,
  lem:minimal-two-sphere-area-variation,
  lem:two-sphere-width-regular-continuity,
  lem:two-sphere-width-surgery-control}
\notready
Along a component path with nonzero $\pi_2$, after all essential sphere
surgeries have ceased, $W_2$ is continuous at regular times, satisfies
\[
 D^+W_2(t)\le-4\pi-\frac12R_{\min}(t)W_2(t)
\]
there, and obeys the left lower-semicontinuity inequality of
\ref{lem:two-sphere-width-surgery-control} at every surgery time.
\end{lemma}

\begin{proof}
The two continuity assertions are
\ref{lem:two-sphere-width-regular-continuity} and
\ref{lem:two-sphere-width-surgery-control}.  At a regular time $t$, choose a
minimizer $f_t$ from \ref{thm:two-sphere-width-realization}.  Keeping this
map fixed for $h>0$ gives
\[
 W_2(t+h)-W_2(t)
 \le\operatorname{Area}_{g(t+h)}(f_t)-\operatorname{Area}_{g(t)}(f_t).
\]
Divide by $h$, take the upper limit as $h\downarrow0$, and apply
\ref{lem:minimal-two-sphere-area-variation}.
\end{proof}

\begin{theorem}[Finiteness of essential sphere surgeries]
\label{thm:finite-essential-sphere-surgeries}
\dcref{MT:ch18:18.7,MT:ch18:18.5}
\uses{def:controlled-ricci-flow-with-surgery}
\notready
Assume that each essential sphere surgery decreases a nonnegative integer
topological complexity by at least one, while the initial complexity is
finite.  Then, in a controlled Ricci flow with surgery, only finitely many
surgery times use
a sphere which is non-null-homotopic in its presurgery component.
\end{theorem}

\begin{proof}
If the initial complexity is $N$, the assumed strict decrease leaves a
nonnegative integer at most $N-j$ after $j$ essential surgeries.  Hence
$j\le N$, proving finiteness.
\end{proof}

\section{Eventual Vanishing of \texorpdfstring{$\pi_2$}{pi2}}

\begin{lemma}[Scalar-curvature lower bound]
\label{lem:extinction-scalar-curvature-lower-bound}
\dcref{MT:ch18:18.1,MT:ch15:15.9}
\uses{thm:all-time-controlled-surgery-flow}
\notready
For the normalized all-time controlled surgery flow, including the
postsurgery time slices,
\[
 R_{\min}(t)\ge-\frac6{1+4t}.
\]
\end{lemma}

\begin{proof}
The pinching-toward-positive conclusion in
\ref{thm:all-time-controlled-surgery-flow} includes the scalar inequality at
regular times.  At a surgery time the inserted cap and retained neck satisfy
the same pinching inequality, while discarded components do not affect the
minimum on a surviving component.  Thus the estimate passes through every
surgery and holds for all time slices.
\end{proof}

\begin{theorem}[Noncompact two-connected covers are contractible]
\label{thm:noncompact-two-connected-cover-contractible}
\dcref{MT:ch18:18.9,Hatcher:ch3:3.29,Hatcher:ch4:4.5,Hatcher:ch4:4.32}
\notready
Let $X$ be a connected three-manifold with $\pi_2(X)=0$.  If its universal
cover $\widetilde X$ is noncompact, then $\widetilde X$ is contractible.
\end{theorem}

\begin{proof}
The cover is simply connected and covering invariance gives
$\pi_2(\widetilde X)=0$.  Degree-two Hurewicz therefore gives
$H_2(\widetilde X;\mathbb Z)=0$, while simple connectivity gives $H_1=0$.
The cover is orientable, and the top homology of a connected noncompact
three-manifold vanishes, so $H_3=0$.  A smooth triangulation has dimension
three, hence $H_i=0$ for $i>3$.

Induct on $n\ge3$.  If $\pi_i(\widetilde X)=0$ for $i<n$, the Hurewicz
theorem identifies $\pi_n(\widetilde X)$ with $H_n(\widetilde X;\mathbb Z)$,
which is zero.  Thus all homotopy groups vanish.  Smooth manifolds and their
covering spaces have CW type, so the map from a point to $\widetilde X$ is a
weak homotopy equivalence between connected CW-type spaces.  Whitehead's
theorem makes it a homotopy equivalence, proving contractibility.
\end{proof}

\begin{theorem}[Eventual vanishing of $\pi_2$]
\label{thm:eventual-pi2-vanishing}
\dcref{MT:ch18:18.9}
\uses{thm:finite-essential-sphere-surgeries,
  lem:two-sphere-width-forward-inequality,
  lem:extinction-scalar-curvature-lower-bound,
  thm:forward-dini-comparison,
  thm:noncompact-two-connected-cover-contractible}
\notready
For an all-time controlled surgery flow there is $T_1<\infty$ such that every
component of every time slice $M_T$, $T\ge T_1$, has trivial $\pi_2$.  Any
such component with infinite fundamental group has contractible universal
cover.
\end{theorem}

\begin{proof}
By \ref{thm:finite-essential-sphere-surgeries}, choose $T_0$ after the last
essential sphere surgery.  A component of $M_{T_0}$ with nonzero $\pi_2$ has
a unique continuation with nonzero $\pi_2$ until it is discarded: surgery on
null-homotopic spheres leaves exactly one descendant homotopy equivalent to
the predecessor, while every other descendant is a homotopy three-sphere.

Suppose one such path persists for all $t\ge T_0$.  By
\ref{lem:two-sphere-width-forward-inequality} and
\ref{lem:extinction-scalar-curvature-lower-bound}, its positive width obeys
\[
 D^+W_2(t)\le-4\pi+\frac3{1+4t}W_2(t).
\]
Let $w(T_0)=W_2(T_0)$ and
$w'=-4\pi+3w/(1+4t)$.  The integrating factor
$(1+4t)^{-3/4}$ gives
\[
 w(t)=W_2(T_0)\left(\frac{1+4t}{1+4T_0}\right)^{3/4}
 +4\pi(1+4T_0)^{1/4}(1+4t)^{3/4}-4\pi(1+4t).
\]
Continuity at regular times, the surgery liminf inequality, and
\ref{thm:forward-dini-comparison} give $W_2(t)\le w(t)$.  The last term in
$w$ has order $t$, while the first two have order $t^{3/4}$, so $w$ is
eventually negative, contradicting $W_2(t)>0$.

There are finitely many components at $T_0$, hence one finite $T_1$ works for
all their continuations.  Triviality of $\pi_2$ persists under later trivial
sphere surgeries.  The final assertion follows from
\ref{thm:noncompact-two-connected-cover-contractible}, since an infinite
fundamental group makes the universal cover noncompact.
\end{proof}
\section{Late-Component \texorpdfstring{$\pi_3$}{pi3} Topology}

\begin{theorem}[Persistence of the fundamental-group class]
\label{thm:component-fundamental-group-persistence}
\dcref{MT:ch18:18.19,MT:ch15:15.3}
\uses{def:controlled-ricci-flow-with-surgery,
  thm:trivial-surgery-component-map,thm:finite-interval-surgery-bound}
\notready
Let $[0,T]$ be a bounded interval containing only finitely many surgery times.
Suppose the fundamental group of every initial component of a controlled
surgery flow is a free product of finitely many finite groups and infinite
cyclic groups.  Then the same is true for every component of every time slice
up to $T$.
\end{theorem}

\begin{proof}
The time vector field identifies components diffeomorphically on every
regular interval.  At a surgery time, the surgery description writes each
presurgery component as a connected sum of the selected postsurgery
components and the discarded spherical or $S^2$-bundle pieces.  Van Kampen
therefore makes the fundamental group of each postsurgery component a free
factor of the corresponding presurgery group; the discarded factors are
finite or infinite cyclic.  This is compatible with the connected-sum pinch
maps in \ref{thm:trivial-surgery-component-map}.

It remains to check the group-theoretic closure.  Let
$G=G_1*\cdots*G_m*F_r$, where the $G_i$ are finite, and let $H$ be a free
factor of $G$.  Restrict the normal-form Bass--Serre tree of this splitting
to $H$.  Edge stabilizers remain trivial, and every vertex stabilizer is an
intersection $H\cap gG_ig^{-1}$, hence finite; the free part of the quotient
graph supplies infinite cyclic factors.  The group $H$ is a retract of the
finitely generated group $G$, so it is finitely generated.  Normal forms then
show that only finitely many nontrivial vertex stabilizers and finitely many
free generators occur.  Thus $H$ has the asserted form.  Induction over the
surgery times proves the theorem.
\end{proof}

\begin{theorem}[Free-product obstruction to vanishing $\pi_2$]
\label{thm:free-product-pi2-obstruction}
\dcref{MT:ch18:18.20}
\notready
Assume the closed three-manifold sphere/free-product theorem in the form that
a nontrivial free product or an infinite cyclic fundamental group yields an
essential (non-null-homotopic) embedded $2$-sphere.  If $X$ is a closed
connected three-manifold with $\pi_1(X)$ a nontrivial free product or
$\pi_1(X)\cong\mathbb Z$,
then
$\pi_2(X)\ne0$.
\end{theorem}

\begin{proof}
The assumed sphere/free-product theorem supplies an embedded sphere whose
inclusion is not null-homotopic.  Its homotopy class is therefore a nonzero
element of $\pi_2(X)$.
\end{proof}

\begin{theorem}[Finite fundamental group and the universal cover]
\label{thm:finite-cover-homotopy-three-sphere}
\dcref{MT:ch18:18.19,Hatcher:ch3:3.30,Hatcher:ch4:4.32}
\notready
Let $X$ be a closed connected three-manifold with finite fundamental group
and $\pi_2(X)=0$.  Its universal cover $\widetilde X$ is a homotopy
three-sphere, and covering projection induces
\[
 \pi_3(\widetilde X)\cong\pi_3(X)\cong\mathbb Z.
\]
\end{theorem}

\begin{proof}
The universal cover is finite-sheeted, hence compact, closed, and simply
connected.  It is orientable, and covering invariance gives
$\pi_2(\widetilde X)=0$.  Poincare duality gives
$H_3(\widetilde X;\mathbb Z)\cong\mathbb Z$, while degree-two Hurewicz gives
$H_2=0$ and simple connectivity gives $H_1=0$.  Degree-three Hurewicz now
identifies $\pi_3(\widetilde X)$ with $H_3$.

Choose a map $S^3\to\widetilde X$ representing a generator.  It is an
isomorphism on every integral homology group.  Both spaces are simply
connected CW-type spaces, so the homological Whitehead theorem makes this map
a homotopy equivalence.  Thus $\widetilde X$ is a homotopy three-sphere.
Covering invariance in degrees at least two gives the asserted isomorphism
with $\pi_3(X)$.
\end{proof}

\begin{theorem}[Nonzero $\pi_3$ on a late surviving component]
\label{thm:late-component-nontrivial-pi3}
\dcref{MT:ch18:18.19}
\uses{thm:eventual-pi2-vanishing,
  thm:component-fundamental-group-persistence,
  thm:free-product-pi2-obstruction,
  thm:finite-cover-homotopy-three-sphere}
\notready
Let $T_1$ be as in \ref{thm:eventual-pi2-vanishing}, and assume the initial
fundamental groups are free products of finite groups and infinite cyclic
groups.  If $\mathcal X$ is a component path defined on $[T_1,T_2]$, then
\[
 \pi_3(\mathcal X(T_1))\cong\mathbb Z.
\]
In particular it contains a nonzero class.
\end{theorem}

\begin{proof}
Put $X_1=\mathcal X(T_1)$.  By \ref{thm:eventual-pi2-vanishing},
$\pi_2(X_1)=0$.  The persistence theorem
\ref{thm:component-fundamental-group-persistence} expresses $\pi_1(X_1)$
as a free product of finite and infinite cyclic
groups.  If more than one nontrivial factor occurred, or if an infinite
cyclic factor occurred, \ref{thm:free-product-pi2-obstruction} would
contradict the vanishing of $\pi_2$.  Thus the fundamental group is finite,
possibly trivial.  Apply \ref{thm:finite-cover-homotopy-three-sphere}.
\end{proof}

\section{The Loop-Space Width Interface}

Let $\Lambda_0X$ denote the component of null-homotopic loops in the free
$C^1$ loop space of a compact manifold $X$, and let $*$ be the constant loop
at a chosen base point $x_0$.

\begin{lemma}[The quotient and wedge model]
\label{lem:loop-space-quotient-wedge}
\dcref{MT:ch18:18.16}
\notready
For $c_0\in S^2$, the quotient
\[
 Q=(S^2\times S^1)/(\{c_0\}\times S^1)
\]
is homotopy equivalent to $S^2\vee S^3$.  Consequently, for every based
space $X$ of CW type,
\[
 [Q,X]_*\cong\pi_2(X)\times\pi_3(X).
\]
\end{lemma}

\begin{proof}
Give $S^2$ one zero-cell and one two-cell, choosing the zero-cell to be
$c_0$, and give $S^1$ one zero-cell and one one-cell.  The product has cells
in dimensions $0,1,2,3$.  Collapsing $\{c_0\}\times S^1$ removes the
one-cell.  The attaching map of the remaining three-cell is the Whitehead
product of the two-cell with the collapsed one-cell, hence is null after the
collapse.  The quotient CW complex is therefore a two-sphere with a
three-cell attached at the base point, namely $S^2\vee S^3$.  Based maps out
of a wedge are pairs of based maps, which gives the final identification.
\end{proof}

\begin{lemma}[Lifting null-loop families]
\label{lem:loop-space-universal-cover-lift}
\dcref{MT:ch18:18.16}
\notready
Let $p\colon\widetilde X\to X$ be the universal cover of a connected
manifold.  A continuous family $\Gamma\colon S^2\to\Lambda_0X$ lifts, after
choosing one lift of one marked point, to a continuous family of loops in
$\widetilde X$.  Homotopies lift in the same way, and changing the initial
lift acts by one deck transformation on the entire lifted family.
\end{lemma}

\begin{proof}
The adjoint map $F\colon S^2\times S^1\to X$ sends every circle
$\{c\}\times S^1$ to a null-homotopic loop.  The induced homomorphism on
$\pi_1(S^2\times S^1)\cong\mathbb Z$ is therefore zero, so the covering
lifting criterion gives a lift $\widetilde F$.  Uniqueness after fixing one
point proves continuity of the family and the assertion about deck
transformations.  Apply the same argument to
$S^2\times S^1\times[0,1]$ for a homotopy.
\end{proof}

\begin{theorem}[Free-loop-space identification of $\pi_3$]
\label{thm:loop-space-pi3-identification}
\dcref{MT:ch18:18.16}
\uses{lem:loop-space-quotient-wedge,
  lem:loop-space-universal-cover-lift}
\notready
If $X$ is compact, connected, and $\pi_2(X,x_0)=0$, then adjunction induces
an isomorphism
\[
 \pi_2(\Lambda_0X,*)\cong\pi_3(X,x_0).
\]
Free homotopy classes of not necessarily based maps
$S^2\to\Lambda_0X$ correspond to the orbit set of $\pi_3(X,x_0)$ under its
natural $\pi_1(X,x_0)$-action.
\end{theorem}

\begin{proof}
For based families, compact-open adjunction identifies a map
$S^2\to\Lambda_0X$ with a map $S^2\times S^1\to X$ which collapses
$\{c_0\}\times S^1$.  By \ref{lem:loop-space-quotient-wedge}, its based
homotopy class is a pair in $\pi_2(X)\times\pi_3(X)$.  The first factor
vanishes, giving the displayed isomorphism.  Passing between continuous and
$C^1$ loops does not change the homotopy type: smoothing in finitely many
charts, relative to already smooth loops, gives mutually inverse homotopies.

For an unbased family in $\Lambda_0X$, use
\ref{lem:loop-space-universal-cover-lift}.  A choice of lift turns it into the
based construction after moving one constant loop to the chosen base point.
A different lift applies a deck transformation throughout and changes the
based representative by the natural $\pi_1(X)$-action.  Conversely, a change
by this action is realized by moving the base point around the corresponding
loop.  Thus forgetting the base point gives precisely the stated orbit set.
\end{proof}

\begin{definition}[Loop-family width]
\label{def:loop-family-width}
\dcref{MT:ch18:18.17}
\uses{thm:loop-space-pi3-identification}
\notready
For a null-homotopic $C^1$ loop $\gamma$ in a compact Riemannian manifold,
let
\[
 \mathcal A_g(\gamma)=\inf_F\operatorname{Area}_g(F),
\]
where $F\colon D^2\to X$ ranges over Lipschitz maps whose boundary agrees
with $\gamma$ up to orientation-preserving reparameterization.  For a
continuous family $\Gamma\colon S^2\to\Lambda_0X$, put
\[
 W_g(\Gamma)=\max_{c\in S^2}\mathcal A_g(\Gamma(c)).
\]
If $\xi\in\pi_3(X)$, let $[\xi]$ be its $\pi_1(X)$-orbit.  By
\ref{thm:loop-space-pi3-identification}, this orbit corresponds to a free
homotopy class of loop families.  Define
\[
 W_g(\xi)=\inf_{[\Gamma]=[\xi]}W_g(\Gamma),
\]
where representatives need not preserve the base point.  Thus $W_g$ is
constant on every $\pi_1(X)$-orbit.
\end{definition}

\begin{lemma}[Continuity of least spanning area]
\label{lem:least-spanning-area-continuity}
\dcref{MT:ch18:18.17}
\notready
The function $\gamma\mapsto\mathcal A_g(\gamma)$ is finite and continuous
on $\Lambda_0X$ in the $C^1$ topology.  More precisely, two sufficiently
$C^1$-close loops cobound a Lipschitz annulus whose area tends to zero with
their $C^1$ distance, and
\[
 |\mathcal A_g(\gamma_0)-\mathcal A_g(\gamma_1)|
 \le\inf\{\operatorname{Area}_g(A):\partial A=\gamma_0\sqcup\gamma_1\}.
\]
\end{lemma}

\begin{proof}
A null-homotopy can be approximated by a Lipschitz disk, so the infimum is
finite.  If $\gamma_0$ and $\gamma_1$ are sufficiently $C^1$-close, join
$\gamma_0(s)$ to $\gamma_1(s)$ by the unique short geodesic.  Smooth
dependence of these geodesics produces an annulus whose two derivatives are
uniformly bounded respectively by the pointwise distance and by the common
$C^1$ bound.  Its area therefore tends to zero with the $C^1$ distance.
Gluing this annulus to a disk spanning either boundary gives both directions
of the displayed inequality and hence continuity.
\end{proof}

\section{The Width Comparison ODE}

\begin{definition}[Comparison equation]
\label{def:loop-width-comparison-ode}
\dcref{MT:ch18:18.23}
\notready
For a compact Ricci flow $(M,g(t))$, $t\in[t_0,t_1]$, and
$t'\in[t_0,t_1]$, let $w_{a,t'}$ be the solution of
\[
 w'=-2\pi-\frac12R_{\min}(t)w,
 \qquad w(t')=a.
\]
Write $w_a=w_{a,t_0}$.
\end{definition}

\begin{lemma}[Integrating-factor formula]
\label{lem:loop-width-ode-formula}
\dcref{MT:ch18:18.25,MT:ch18:18.26}
\uses{def:loop-width-comparison-ode}
\notready
Put
\[
 B_{t'}(t)=\int_{t'}^t\frac12R_{\min}(s)\,ds.
\]
Then, for $t''\ge t'$,
\[
 w_{a,t'}(t'')=e^{-B_{t'}(t'')}
 \left(a-2\pi\int_{t'}^{t''}e^{B_{t'}(s)}\,ds\right).
\]
If $a'\ge a$, then
\[
 w_{a',t'}(t'')-w_{a,t'}(t'')
 =(a'-a)e^{-B_{t'}(t'')}\ge0.
\]
On a fixed compact time interval this difference is bounded by a constant
times $a'-a$.
\end{lemma}

\begin{proof}
Multiplication of the equation in \ref{def:loop-width-comparison-ode} by
$e^{B_{t'}(t)}$ gives
\[
 \frac d{dt}\bigl(e^{B_{t'}(t)}w_{a,t'}(t)\bigr)
 =-2\pi e^{B_{t'}(t)}.
\]
Integration from $t'$ to $t''$ proves the first formula.  Subtracting the
formulas with initial values $a'$ and $a$ proves the second.  Boundedness of
$R_{\min}$ on a compact regular interval bounds the exponential factor.
\end{proof}
\section{Controlled Loop-Family Deformation}

\begin{theorem}[Least-disk variation interface]
\label{thm:least-disk-area-variation-interface}
\dcref{MT:ch19:19.2,MT:ch19:19.4}
\uses{def:loop-family-width}
\notready
Assume Plateau minimizers for the relevant null-homotopic curves exist with
the required boundary regularity, and that their moving-boundary first
variation is valid.  Let $g(t)$ be a Ricci flow on a compact three-manifold and let
$\gamma(t)$ be a smooth curve-shrinking flow of null-homotopic immersed
loops.  Then $t\mapsto\mathcal A_{g(t)}(\gamma(t))$ is continuous and, at
every time at which the flow is smooth,
\[
 D^+\mathcal A_{g(t)}(\gamma(t))
 \le-2\pi-\frac12R_{\min}(t)
       \mathcal A_{g(t)}(\gamma(t)).
\]
\end{theorem}

\begin{theorem}[Regularity of an attained annular minimizer]
\label{thm:annular-douglas-minimizer}
\dcref{MT:ch19:19.15,
  white-classical-area-minimizing-surfaces:page-0001,
  white-classical-area-minimizing-surfaces:page-0002}
\notready
Let $(N,h)$ be a compact real-analytic Riemannian manifold of dimension at
least three.  Let $c_0,c_1$ be disjoint embedded real-analytic closed curves
which are freely homotopic and homotopically nontrivial.  Suppose an admissible
spanning annulus $F_0$ attains the annular area infimum, and that there is
$\rho>0$ such that, for every essential simple closed curve $\alpha$ in its
domain,
\[
 L_h(F_0\circ\alpha)\ge\rho.
\]
After replacing $F_0$ by its conformal harmonic representative in the same
minimizing class, it has finite interior branch set, is real analytic up to
its real-analytic boundary, and has no true boundary branch point.  False
boundary branch points are not excluded.  There are interior parallel
boundary curves avoiding the finite branch set which converge to the original
boundary and whose intervening collar areas tend to zero.
\end{theorem}

\begin{lemma}[Ramp nondegeneration]
\label{lem:ramp-annulus-nondegeneration}
\dcref{MT:ch19:19.15}
\uses{def:ramp-curve}
\notready
Let $c_0,c_1\colon S^1\to M\times S^1_\lambda$ be degree-one ramps in the
same free homotopy class, and let
$F\colon S^1\times[0,1]\to M\times S^1_\lambda$ span them.  For every
essential simple closed domain curve $\alpha$,
\[
 L(F\circ\alpha)\ge L(S^1_\lambda)>0.
\]
The same bound holds while both curves evolve by smooth curve shrinking.
\end{lemma}

\begin{proof}
An essential simple closed curve in the annulus is freely homotopic, with one
of the two orientations, to either boundary component.  Its composition with
the circle projection therefore has degree $1$ or $-1$.  Such a map onto
$S^1_\lambda$ has length at least the circumference of the target, and product
projection does not increase length.  A smooth curve flow is a homotopy, so
the circle degree and common free homotopy class are preserved.
\end{proof}

\begin{lemma}[Variation of the annular area infimum]
\label{lem:annular-area-infimum-variation}
\dcref{MT:ch19:19.15}
\uses{def:ramp-curve,thm:annular-douglas-minimizer,
  lem:ramp-annulus-nondegeneration}
\notready
Let $n=\dim M\ge3$, and let $c_0(t),c_1(t)$ be degree-one ramp
curve-shrinking flows in $(M,g(t))\times S^1_\lambda$ in the same free
homotopy class.  Assume that, for each time under consideration, the annular
infimum is attained by a minimizer satisfying the hypotheses of
\ref{thm:annular-douglas-minimizer}.  If $\mu(t)$ is the infimum of the areas
of spanning annuli,
then $\mu$ is continuous and
\[
 D^+\mu(t)\le(2n-1)\max_M|\operatorname{Rm}_{g(t)}|\,\mu(t).
\]
Consequently $\mu(t)\le e^{C(t-s)}\mu(s)$ on a compact time interval, with
$C$ independent of $\lambda$.
\end{lemma}

\begin{proof}
First suppose the boundary curves are disjoint and real analytic.  By
\ref{lem:ramp-annulus-nondegeneration}, the hypothesis of
\ref{thm:annular-douglas-minimizer} holds.  Choose a minimizing annulus and
move its boundary by the two curve flows.  Interior normal variation vanishes
by minimality.  Metric and boundary variation give
\[
 \left.\frac d{ds}\right|_{s=t}\operatorname{Area}(A_s)
 =-\int_A\operatorname{tr}_{TA}\operatorname{Ric}\,da
   -\int_{\partial A}k_g\,ds.
\]
Apply Gauss--Bonnet on parallel subannuli avoiding the finite branch set and
let them tend to the boundary.  If the interior branch orders are $m_j$, the
result is
\[
 -\int_{\partial A}k_g\,ds
 =\int_AK_A\,da-2\pi\sum_j(m_j-1)
 \le\int_AK_A\,da.
\]
Minimality gives $K_A\le\sec_{M\times S^1_\lambda}(TA)$, and
$|\operatorname{tr}_{TA}\operatorname{Ric}|\le2(n-1)\max|\operatorname{Rm}|$.  The circle is flat, so
the constants are independent of $\lambda$.  This proves the upper
derivative estimate.

In dimension at least three, perturb one boundary curve in $C^2$ so that the
curves are disjoint; the original and perturbed curves cobound a thin annulus
whose area tends to zero.  Gluing this annulus to competitors compares the
two annular infima in both directions.  The same gluing argument using the
thin annuli swept out over $[t,t+h]$ proves upper and lower semicontinuity.
Letting the perturbation vanish proves continuity and the differential
inequality in general.  Scalar integrating-factor comparison gives the final
exponential estimate.
\end{proof}

\begin{theorem}[Small-annulus length transfer]
\label{thm:small-annulus-length-transfer}
\dcref{MT:ch19:19.31,MT:ch19:19.35}
\uses{def:ramp-curve,thm:annular-douglas-minimizer,
  lem:ramp-annulus-nondegeneration}
\notready
Assume the local quantitative annular graph estimate for ramps: under the
curvature and small-turning hypotheses displayed below, an annulus of area
less than $\mu_0$ cannot reduce the boundary length by more than the stated
factor, uniformly in $\lambda$.  There is $\delta_0>0$ with the following
property.  For every $r>0$ and every
ambient sectional-curvature upper bound $K$, there is $\mu_0>0$ such that if
two ramps $\gamma,\widehat\gamma$ in $M\times S^1_\lambda$ cobound an
annulus of area below $\mu_0$, and
\[
 L(\gamma)\ge r,
 \qquad
 \int_I|H_\gamma|\,ds<\delta_0
\]
for every length-$r$ subarc $I\subset\gamma$, then
\[
 L(\widehat\gamma)\ge\frac34L(\gamma).
\]
The constants are uniform in $\lambda$.
\end{theorem}

\begin{proof}
This is precisely the local quantitative annular graph estimate assumed in
the statement, with its constants chosen uniformly in $\lambda$.
\end{proof}

\begin{theorem}[Controlled deformation of a loop family]
\label{thm:controlled-loop-family-deformation}
\dcref{MT:ch18:18.24}
\uses{def:loop-family-width,def:loop-width-comparison-ode,
  lem:loop-width-ode-formula,lem:least-spanning-area-continuity,
  thm:least-disk-area-variation-interface,
  lem:annular-area-infimum-variation,
  thm:small-annulus-length-transfer,
  thm:controlled-curve-shrinking-existence}
\notready
Assume the least-disk, attained-annulus, annular-transfer, and compact-family
curve-shrinking interfaces in the dependency list hold uniformly on
$[t_0,t_1]$, including their stated regularity and nondegeneracy hypotheses.
Let $g(t)$ be a Ricci flow on a compact three-manifold for
$t\in[t_0,t_1]$, let $\Gamma\colon S^2\to\Lambda_0M$ be continuous, and
let $\zeta>0$.  There is a continuous family
$\widetilde\Gamma(t)\colon S^2\to\Lambda_0M$ such that
$[\widetilde\Gamma(t_0)]=[\Gamma]$ and, for every $c\in S^2$,
\[
 |\mathcal A_{g(t_0)}(\widetilde\Gamma(t_0)(c))
   -\mathcal A_{g(t_0)}(\Gamma(c))|<\zeta,
\]
and at least one of the following holds:
\begin{enumerate}
\item $L_{g(t_1)}(\widetilde\Gamma(t_1)(c))<\zeta$;
\item
\[
 \mathcal A_{g(t_1)}(\widetilde\Gamma(t_1)(c))
 \le w_{\mathcal A_{g(t_0)}(\widetilde\Gamma(t_0)(c)),t_0}(t_1)+\zeta.
\]
\end{enumerate}
\end{theorem}

\section{\texorpdfstring{$\pi_3$}{pi3}-Width Evolution Through Surgery}

\begin{theorem}[Transport of a nonzero third-homotopy class]
\label{thm:surgery-pi3-class-transport}
\dcref{MT:ch18:18.18,MT:ch15:15.12}
\uses{thm:trivial-surgery-component-map}
\notready
Let $\mathcal X$ be a component path on $[t_0,t_1]$ with only finitely many
surgery times, with
$\pi_2(\mathcal X(t_0))=0$ and
$\pi_3(\mathcal X(t_0))\cong\mathbb Z$, and let
$0\ne\xi(t_0)\in\pi_3(\mathcal X(t_0))$.  Then $\pi_2(\mathcal X(t))=0$
for every $t$, and there is a compatible family of nonzero classes
$\xi(t)\in\pi_3(\mathcal X(t))$.  At a surgery time $T$, the maps
$F_t\colon\mathcal X(t)\to\mathcal X(T)$ of
\ref{thm:trivial-surgery-component-map}, for $t<T$ sufficiently close to $T$,
satisfy
\[
 (F_t)_*\xi(t)=\xi(T).
\]
\end{theorem}

\begin{proof}
At a regular time, the flow identification transports the homotopy groups and
the class.  At a surgery time, every cutting sphere is null homotopic because
$\pi_2=0$.  Apply \ref{thm:trivial-surgery-component-map}.  It preserves the
vanishing of $\pi_2$, and its map on the relevant infinite cyclic third
homotopy groups is injective.  Hence
$\xi(T)=(F_t)_*\xi(t)$ is nonzero.  Induction over the surgery times gives the
compatible family.
\end{proof}

For the transported class in \ref{thm:surgery-pi3-class-transport}, write
$W_\xi(t)=W_{g(t)}(\xi(t))$.

\begin{lemma}[Continuity of loop-family width at regular times]
\label{lem:loop-width-regular-continuity}
\dcref{MT:ch18:18.21}
\uses{def:loop-family-width,thm:surgery-pi3-class-transport}
\notready
The function $W_\xi$ is continuous at every nonsurgery time.
\end{lemma}

\begin{proof}
Use the time vector field to identify a closed surgery-free neighborhood of
$t$ with a fixed compact manifold.  The pulled-back metrics satisfy
\[
 c(s)^{-1}g(t)\le g(s)\le c(s)g(t),
 \qquad c(s)\longrightarrow1.
\]
Every Lipschitz disk then has its area multiplied by a factor between
$c(s)^{-1}$ and $c(s)$.  The same bounds pass first to the infimum defining
$\mathcal A$, then to the maximum over a family, and finally to the infimum
over representatives of the transported class.  Thus
\[
 c(s)^{-1}W_\xi(t)\le W_\xi(s)\le c(s)W_\xi(t),
\]
which proves continuity.
\end{proof}

\begin{lemma}[Loop-family width at a surgery time]
\label{lem:loop-width-surgery-lower-semicontinuity}
\dcref{MT:ch18:18.22,MT:ch15:15.12}
\uses{def:loop-family-width,thm:surgery-pi3-class-transport,
  thm:trivial-surgery-component-map}
\notready
At a surgery time $T$,
\[
 W_\xi(T)\le\liminf_{t\uparrow T}W_\xi(t).
\]
\end{lemma}

\begin{proof}
For every $\eta>0$ and all $t<T$ sufficiently close to $T$,
\ref{thm:trivial-surgery-component-map} gives a map $F_t$ with Lipschitz
constant at most $1+\eta$.  By
\ref{thm:surgery-pi3-class-transport}, it sends $\xi(t)$ to $\xi(T)$.  If a
Lipschitz disk spans $\gamma$, then $F_t$ composed with that disk spans
$F_t\circ\gamma$ and has area at most $(1+\eta)^2$ times the original area.
Apply this pointwise to a representative family and take both infima to get
\[
 W_\xi(T)\le(1+\eta)^2W_\xi(t).
\]
Take a sequence realizing the left liminf and then let $\eta\downarrow0$.
\end{proof}

\begin{lemma}[Short loops bound small-area disks]
\label{lem:short-loops-small-disk}
\dcref{MT:ch18:18.28}
\notready
Let $(X,g)$ be compact.  For every $\eta>0$ there is
$0<\zeta<\eta/2$ such that every $C^1$ loop of length less than $\zeta$
bounds a Lipschitz disk of area less than $\eta$.
\end{lemma}

\begin{proof}
Choose $r>0$ below the normal radius uniformly over $X$.  If a loop $\gamma$
has length $\ell<r/2$, put $p=\gamma(1)$ and write
$v(\theta)=\exp_p^{-1}(\gamma(\theta))$.  The radial filling
\[
 F(s,\theta)=\exp_p(sv(\theta))
\]
is defined on the disk.  On the compact radius-$r$ normal bundle the
differential of the exponential map is uniformly bounded.  Hence
$|\partial_sF|\le C\ell$ and
$|\partial_\theta F|\le C|v'(\theta)|$, so
\[
 \operatorname{Area}(F)\le C^2\ell\int_{S^1}|v'|\,d\theta\le C'\ell^2.
\]
Choose
$\zeta<\min\{\eta/2,r/2,\sqrt{\eta/C'}\}$.  Then every loop of length below
$\zeta$ has the required filling.
\end{proof}

\begin{lemma}[Short loop families are trivial]
\label{lem:short-loop-families-trivial}
\dcref{MT:ch18:18.27}
\uses{thm:loop-space-pi3-identification}
\notready
Let $(X,g)$ be compact with $\pi_2(X)=0$.  There is $\zeta>0$ such that any
continuous family $\Gamma\colon S^2\to\Lambda_0X$ all of whose loops have
length less than $\zeta$ represents the zero free-homotopy orbit, and hence
the zero class after any based choice, under
$\pi_2(\Lambda_0X,*)\cong\pi_3(X)$.
\end{lemma}

\begin{proof}
Choose $\zeta$ below the normal radius.  For each $c\in S^2$, join
$\Gamma(c)(\theta)$ to $\Gamma(c)(1)$ by the unique short geodesic.  Smooth
dependence of the geodesic on its endpoints gives a continuous family of
radial disks.  Shrinking through those disks homotopes $\Gamma$ to the family
of constant loops $c\mapsto\Gamma(c)(1)$.  This remaining family is a map
$S^2\to X$, and is null-homotopic because $\pi_2(X)=0$.  The identification
in \ref{thm:loop-space-pi3-identification} gives the conclusion.
\end{proof}

\begin{theorem}[Forward inequality for the $\pi_3$ width]
\label{thm:loop-width-forward-inequality}
\dcref{MT:ch18:18.18}
\uses{def:loop-family-width,def:loop-width-comparison-ode,
  lem:loop-width-ode-formula,thm:controlled-loop-family-deformation,
  lem:short-loops-small-disk,lem:short-loop-families-trivial,
  lem:loop-width-regular-continuity,
  lem:loop-width-surgery-lower-semicontinuity,
  thm:surgery-pi3-class-transport}
\notready
Let $\mathcal X$ be a component path of a controlled surgery flow, suppose
$\pi_2(\mathcal X(t_0))=0$ and
$\pi_3(\mathcal X(t_0))\cong\mathbb Z$, and let
$0\ne\xi(t_0)\in\pi_3(\mathcal X(t_0))$.  For the transported class,
$W_\xi$ is continuous at nonsurgery times, obeys
\[
 W_\xi(T)\le\liminf_{t\uparrow T}W_\xi(t)
\]
at surgery times, and satisfies
\[
 D^+W_\xi(t)
 \le-2\pi-\frac12R_{\min}(t)W_\xi(t)
\]
on the immediate right of every time in the component path.
\end{theorem}

\begin{proof}
The transport and continuity assertions are
\ref{thm:surgery-pi3-class-transport},
\ref{lem:loop-width-regular-continuity}, and
\ref{lem:loop-width-surgery-lower-semicontinuity}.  Fix a surgery-free
interval $[s,t]$ and let $w$ solve
\ref{def:loop-width-comparison-ode} with $w(s)=W_\xi(s)$.  We prove
$W_\xi(t)\le w(t)$ whenever the component survives to $t$.

First suppose $w(t)\ge0$.  Given $\eta>0$, choose a representative family
with width below $W_\xi(s)+\varepsilon$.  Use
\ref{lem:loop-width-ode-formula} and then choose the deformation error
$\zeta$ so small that changing any initial disk area by at most
$2\zeta$ changes its terminal ODE value by less than $\eta/2$.  Also choose
$\zeta$ by \ref{lem:short-loops-small-disk} so that every terminal loop of
length below $\zeta$ spans a disk of area below $\eta$.  Apply
\ref{thm:controlled-loop-family-deformation}.  In its area alternative the
terminal area is below $w(t)+\eta$; in its short alternative it is below
$\eta\le w(t)+\eta$.  Hence the terminal family has width at most
$w(t)+\eta$.  Let $\eta,\varepsilon\downarrow0$.

If $w(t)<0$, choose $\eta>0$ with $w(t)+\eta<0$ and choose $\zeta$ small
enough for both the ODE sensitivity estimate and
\ref{lem:short-loop-families-trivial}.  The area alternative in
\ref{thm:controlled-loop-family-deformation} would give a negative spanning
area and is impossible.  Thus every terminal loop is short, so the terminal
family is homotopically trivial.  This contradicts transport of the nonzero
class.  Consequently the component cannot survive to a time at which the
comparison solution is negative.

Apply the comparison just proved on arbitrarily short regular intervals.
Together with regular-time continuity it yields the displayed upper forward
derivative.  Starting immediately after a surgery gives the same right-hand
inequality there; the separate left lower-semicontinuity statement accounts
for the jump from the preceding interval.
\end{proof}

\section{Finite Extinction}

\begin{theorem}[The free-product hypothesis excludes a two-sided projective
plane]
\label{thm:free-product-excludes-two-sided-projective-plane}
\dcref{MT:ch18:18.1}
\notready
Assume the corresponding closed three-manifold projective-plane theorem: a
two-sided embedded $\mathbb{RP}^2$ gives a forbidden free-product factor in
the fundamental group.  Let $M$ be a closed connected three-manifold whose
fundamental group is a
free product of finite groups and infinite cyclic groups.  Then $M$ contains
no embedded $\mathbb{RP}^2$ with trivial normal bundle.
\end{theorem}

\begin{proof}
The assumed projective-plane group theorem rules out exactly such an
embedding under the displayed free-product hypothesis.
\end{proof}

\begin{theorem}[Finite-time extinction]
\label{thm:finite-time-extinction}
\dcref{MT:ch18:18.1}
\uses{thm:free-product-excludes-two-sided-projective-plane,
  thm:all-time-controlled-surgery-flow,
  thm:eventual-pi2-vanishing,
  thm:late-component-nontrivial-pi3,
  thm:loop-width-forward-inequality,
  lem:extinction-scalar-curvature-lower-bound,
  thm:forward-dini-comparison}
\notready
Let $(M,g_0)$ be a compact connected normalized Riemannian three-manifold
whose fundamental group is a free product of finite groups and infinite
cyclic groups.  The all-time controlled Ricci flow with surgery starting from
$(M,g_0)$ has a finite extinction time: there is $T_*<\infty$ such that
$M_t=\varnothing$ for every $t\ge T_*$.
\end{theorem}

\begin{proof}
By \ref{thm:free-product-excludes-two-sided-projective-plane}, the initial
manifold satisfies the topological hypothesis of
\ref{thm:all-time-controlled-surgery-flow}; take the flow supplied there.
Choose $T_1$ from \ref{thm:eventual-pi2-vanishing}.  If a component path
$\mathcal X$ survives on $[T_1,T]$, then
\ref{thm:late-component-nontrivial-pi3} supplies a nonzero class
$\xi\in\pi_3(\mathcal X(T_1))$.

By \ref{thm:loop-width-forward-inequality} and
\ref{lem:extinction-scalar-curvature-lower-bound},
\[
 D^+W_\xi(t)\le-2\pi+\frac3{1+4t}W_\xi(t).
\]
Let $w(T_1)=W_\xi(T_1)$ and
$w'=-2\pi+3w/(1+4t)$.  Direct integration gives the consistent formula
\[
 w(t)=W_\xi(T_1)
       \left(\frac{1+4t}{1+4T_1}\right)^{3/4}
 +2\pi(1+4T_1)^{1/4}(1+4t)^{3/4}-2\pi(1+4t).
\]
The comparison principle \ref{thm:forward-dini-comparison}, applied between
successive surgery times and combined with the downward surgery inequality,
gives $0\le W_\xi(t)\le w(t)$ while the path survives.  The negative linear
term in $w$ dominates its $t^{3/4}$ terms, so $w$ becomes negative after a
finite time depending only on $T_1$ and $W_\xi(T_1)$.  Thus this component
path must end before that time.

The compact time slice $M_{T_1}$ has finitely many components.  Taking the
maximum of their finite lifetime bounds gives $T_*<\infty$ after which no
component remains.  Once a time slice is empty, every later slice is empty by
the definition of the surgery flow.
\end{proof}
